\documentclass{article}
\usepackage{graphicx} % Required for inserting images

% Added stuff
\usepackage{amsmath}
\usepackage{amssymb}
\newcommand{\N}{\mathbb{N}}
\newcommand{\Z}{\mathbb{Z}}
\newcommand{\Q}{\mathbb{Q}}
\newcommand{\R}{\mathbb{R}}
\newcommand{\C}{\mathbb{C}}
\newcommand{\bigp}[1]{\left(#1\right)}
\newcommand{\bigb}[1]{\left[#1\right]}
\newcommand{\bigabs}[1]{\left|#1\right|}
\newcommand{\summation}[1][0]{\sum_{n=#1}^\infty}
\newcommand{\summationk}[1][0]{\sum_{k=#1}^\infty}
\newcommand{\summationm}[1][0]{\sum_{m=#1}^\infty}
\newcommand{\eps}{\varepsilon}
\renewcommand{\emptyset}{\varnothing}

\title{p-adic Analysis 1}
\author{Siddharth K}
\date{April 2026}

\begin{document}

\maketitle

\section*{Problem 2}
Let $x, y, z \in \R^n$. We wish to show that
$$\|x - z\| \le \|x - y\| + \|y - z\|.$$
If we redefine $u = x - y$ and $v = z - y$, this is equivalent to
$$\|u + v\| \le \|u\| + \|v\|.$$
Squaring the inequality gives us
$$\|u + v\|^2 \le \|u\|^2 + \|v\|^2 + 2\|u\|\|v\|.$$
We have that
$$\|u + v\|^2 = (u + v)\cdot (u + v) = \|u\|^2 + \|v^2\| + 2(u \cdot v),$$
and by the Cauchy-Schwarz inequality,
$$u \cdot v \le \|u\|\|v\|,$$
$$\|u + v\|^2 = \|u\|^2 + \|v^2\| + 2(u \cdot v) \le \|u\|^2 + \|v^2\| + 2\|u\|\|v\|,$$
as desired.

\section*{Problem 3}
In $\R$, let
$$S_n = (-1/n, 1 + 1/n).$$
Then
$$\bigcap_{n \in \N} S_n = [0, 1],$$
which is closed.

\section*{Problem 6}
Let $x, y, z \in \R$. We want prove the triangle inequality:
$$|e^x - e^y| + |e^y - e^z| \le |e^x - e^z|.$$
Letting $a, b, c = e^x, e^y, e^z$, this is equivalent to the regular triangle inequality on the standard metric on $\R$.
This metric space is not complete. Take $x_n = -n$. This is Cauchy since $e^{-n} \to 0$, but there is no $x \in \R$ such that $e^x = 0$. The completion would be $[0, \infty]$, adding $0$ as the only limit of the Cauchy sequences that don't converge.

\section*{Problem 7}
Let $\eps > 0$. There exists $N$ such that if $i > N$ then
$$d(a_i, b_i) < \eps/2, \quad d(b_i, c_i) < \eps/2.$$
By triangle inequality,
$$d(a_i, c_i) \le d(a_i, b_i) + d(b_i, c_i) < 2 \cdot \eps/2 = \eps,$$
so $a$ is equivalent to $c$.

\section*{Problem 8}
Since
$$\lim_{p \to \infty} d_p(x, y) = \lim_{p \to \infty} \bigp{\sum (y_i - x_i)^p}^{1/p} = \max_i(y_i - x_i),$$
$d_\infty(x, y)$ is the maximum difference between any two individual coordinates of $x, y$.

\section*{Problem 9}
Let $x, y \in X$ and $d = \|y - x\|$. Let $U = B_{d/2}(x), V = B_{d/2}(y)$. Then $U \cap V = \emptyset$ because, if a point $z$ is $< d/2$ away from $x$, by triangle inequality, it must be $> d/2$ away from $y$.

\section*{Problem 12}
Assume a Cauchy sequence $a_n$ converges to $x$ and $y$ with $x \neq y$. Let $0 < \eps < d(x, y) / 2$. We have that $B_\eps(x) \cap B_\eps(y) = \emptyset$, but since $a_n$ converges to $x$ and $y$, there exists $N$ such that if $n > N$, $a_n \in$ $B_\eps(x)$ and $B_\eps(y)$, which is a contradiction. So $x = y$.

\section*{Problem 13}
Consider the sequence
$$x_n = \sum_{k = 0}^n p^k.$$
If $m > n$ then
$$x_m - x_n = p^{n + 1}(1 + p + \cdots + p^{m - n - 1},$$
so
$$d_p(x_m - x_n) = p^{-(n + 1)} \to 0 \text{ as } n \to \infty.$$
So $x_n$ is Cauchy in $\Q$, but there is no $x = \sum_{k = 0}^\infty p^k$ in $\Q$, so the metric space is not complete.

\section*{Problem 16}
Assume $f: (X, d_X) \to (Y, d_Y)$ satisfies the $\eps$-$\delta$ definition. Take any open $U \subseteq Y$. We want to show that $f^{-1}(U)$ is open in $X$.

Pick any $x_0 \in f^{-1}(U)$, so that $f(x_0) \in U$. Since $U$ is open, there exists $\eps > 0$ such that
$$B_Y(f(x_0), \eps) \subseteq U.$$
Since $f$ is continuous at $x_0$, there exists $\delta > 0$ such that
$$d_X(x_0, x) < \delta \implies d_Y(f(x_0), f(x)) < \eps.$$
So if $x \in B_X(x_0, \delta)$, then $f(x) \in B_Y(f(x_0), \eps) \subseteq U$, so $x \in f^{-1}(U)$. Therefore, $B_x(x_0, \delta) \in f^{-1}(U)$, and $f^{-1}(U)$ is open.

Now assume that for every open $U \subseteq Y$, $f^{-1}(U)$ is open in $X$. Fix $x_0 \in X$ and $\eps > 0$. Let
$$U = B_Y(f(x_0), \eps).$$
Since $U$ is open, $f^{-1}(U)$ is open in $X$, and also $x_0 \in f^{-1}(U)$. So there exists $\delta > 0$ such that
$$B_X(x_0, \delta) \subseteq f^{-1}(U).$$
Therefore, whenever $d_X(x_0, x) < \delta$, $x \in f^{-1}(U)$, so
$$f(x) \in U = B_Y(f(x_0), \eps)$$
so $d_Y(f(x_0), f(x)) < \eps$.

\end{document}